% =====================================================================
\chapter{Perceptron}
% =====================================================================

\labinfo{07\_Perceptron}{perceptron.py}{UCI Iris: 100 flowers from two linearly separable species, 4
measurements}{Binary classification (setosa or versicolor)}{Test accuracy 100\%, F1 1.000; 5-fold CV accuracy 100\%}

\section{Dataset}
The binary Iris subset contains 50 setosa and 50 versicolor flowers described by sepal length, sepal
width, petal length and petal width. These two species are linearly separable in this feature space, so
the dataset gives the single-neuron perceptron a clean convergence example. I standardised the features
and made a stratified 80/20 split (80 training, 20 test flowers).

\section{Code and experiments}
\begin{itemize}
\item \code{hardlim} is the transfer function.
\item \code{Perceptron.fit} is a direct translation of the learning rule: shuffle, then for each example
      compute $a$ and $e$ and update $W$ and $b$. After each epoch it records the number of training
      errors and keeps the best weights (the pocket). It can also save snapshots of the weights so the
      moving boundary can be drawn.
\item \code{binary\_metrics} computes accuracy, precision, recall, specificity, F1 and the confusion matrix.
\end{itemize}
\noindent Experiments: (1)~training on all four features for up to 100 epochs and testing;
(2)~5-fold cross-validation on the whole dataset; (3)~the XOR problem, as a check that the
implementation shows the known limitation; (4)~a two-feature perceptron (variance and skewness) to draw
how the boundary moves during training.

\section{Results}

\begin{table}[H]
\centering\small
\caption{Perceptron results on the binary Iris data.}
\label{tab:perc}
\begin{tabular}{@{}l l@{}}
\toprule
Epochs run & 1 (zero errors reached)\\
Fewest training errors & 0 of 80, at epoch 1\\
Final weights & $W=[0.772,\;-1.254,\;1.379,\;1.429]$, \ $b=0$\\
Training accuracy & 100\%\\
\midrule
Test accuracy & \textbf{100\%}\\
Test precision / recall (versicolor) & 1.0000 / 1.0000\\
Test specificity / F1 & 1.0000 / 1.0000\\
Test confusion matrix & [[TN 10, FP 0], [FN 0, TP 10]]\\
\midrule
5-fold CV accuracy & $100\% \pm 0\%$ (all folds 100\%)\\
XOR after 100 epochs & still 2 of 4 wrong (never 0)\\
\bottomrule
\end{tabular}
\end{table}

A single neuron with four weights and a bias classifies 98.5\% of unseen notes correctly, and the
cross-validation folds all lie between 98.2\% and 99.3\%, so the result is stable. The error curve shows
the non-separable behaviour that the theory predicts: the rule never stops changing the weights. The
pocket helped a little here. The weights after the very last epoch made 10 training errors (and scored
98.2\% on the test set), while the best set, from epoch 39, made only 7.

All important weights are negative and of similar size for variance, skewness and curtosis, while
entropy has almost no weight. So forged notes are mainly the ones with low wavelet variance, skewness and
curtosis.

\begin{figure}[H]
\centering
\includegraphics[width=0.48\linewidth]{\figpath{07\_Perceptron}{errors_per_epoch.pdf}}\hfill
\includegraphics[width=0.48\linewidth]{\figpath{07\_Perceptron}{boundary_two_features.pdf}}
\caption{Perceptron convergence and the moving two-feature boundary.}
\end{figure}

\section{Source code}
\lstinputlisting[caption={\texttt{07\_Perceptron/perceptron.py}}]{07_Perceptron/perceptron.py}
\clearpage
